% Part: first-order-logic
% Chapter: syntax-and-semantics
% Section: main-operator

\documentclass[../../../include/open-logic-section]{subfiles}

\begin{document}

\olfileid[pt]{fol}{syn}{mai}

\olsection{\printtoken{S}{main operator} de uma Fórmula}

\begin{explain}
Frequentemente é útil falar sobre o último operador usado na
construção de !!a{formula}~$!A$.  Esse operador é chamado de \emph{operador
  principal} de~$!A$. Intuitivamente, é o operador ``mais externo''
de $!A$. Por exemplo, o operador principal de $\lnot !A$ é $\lnot$,
o operador principal de $(!A \lor !B)$ é $\lor$, etc.
\end{explain}


\begin{defn}[!!^{main operator}]
\ollabel{def:main-op}
O \emph{!!{main operator}} de !!a{formula}~$!A$ é
definido como segue:
\begin{enumerate}
\item \indcase*{!A}{!A}{$\indfrm$ não tem !!{main operator}.}

\tagitem{prvNot}{\indcase{!A}{\lnot !B}{o !!{main operator} de $\indfrm$
  é~$\lnot$.}}{}

\tagitem{prvAnd}{\indcase{!A}{(!B \land !C)}{o !!{main operator} de
  $\indfrm$ é~$\land$.}}{}

\tagitem{prvOr}{\indcase{!A}{(!B \lor !C)}{o !!{main operator} de
  $\indfrm$ é~$\lor$.}}{}

\tagitem{prvIf}{\indcase{!A}{(!B \lif !C)}{o !!{main operator} de
  $\indfrm$ é~$\lif$.}}{}

\tagitem{prvIff}{\indcase{!A}{(!B \liff !C)}{o !!{main operator} de
  $\indfrm$ é~$\liff$.}}{}

\tagitem{prvAll}{\indcase{!A}{\lforall[x][!B]}{o !!{main operator}
  de $\indfrm$ é~$\lforall$.}}{}

\tagitem{prvEx}{\indcase{!A}{\lexists[x][!B]}{o !!{main operator} de
  $\indfrm$ é~$\lexists$.}}{}
\end{enumerate}
\end{defn}

Em cada caso, pretendemos a \emph{ocorrência} específica indicada do
!!{main operator} na fórmula. Por exemplo, como a fórmula
$((!D \lif !E) \lif (!E \lif !D))$ é da forma $(!B \lif !C)$ onde
$!B$ é $(!D \lif !E)$ e $!C$ é $(!E \lif !D)$, a segunda
ocorrência de $\lif$ é o !!{main operator}.

\begin{explain}
Esta é uma definição \emph{recursiva} de uma função que mapeia todas as
!!{formula}s não atômicas para sua ocorrência de !!{main operator}. Por causa
da maneira como as !!{formula}s são definidas indutivamente, toda
!!{formula}~$!A$ satisfaz um dos casos em \olref{def:main-op}.
Isso garante que, para cada !!{formula}~$!A$ não atômica, existe um !!{main
  operator}. Como cada !!{formula} satisfaz apenas uma dessas
condições, e como as !!{formula}s menores a partir das quais $!A$ é
construída são determinadas de forma única em cada caso, a ocorrência de
!!{main operator} de~$!A$ é única, e assim definimos uma
função.
\end{explain}

Chamamos as !!{formula}s pelos nomes em \olref{tab:main-op} dependendo de
qual símbolo é seu !!{main operator}.%
\iftag{defNot,defOr,defAnd,defIf,defIff,defTrue,defFalse,defEx,defAll}
{ Lembre-se, no entanto, de que os operadores definidos não aparecem oficialmente em
!!{formula}s. Eles são apenas abreviações, então oficialmente não podem
ser o operador principal de uma fórmula. Em provas sobre todas as !!{formula}s,
eles, portanto, não precisam ser tratados separadamente.}

\begin{table}[!h]
\centering
\begin{tabular}{c | c | c}
!!^{main operator} & Tipo de !!{formula} & Exemplo\\
\hline
nenhum & (!!{formula}) atômica &
\iftag{prvFalse}{$\lfalse$,}{}
\iftag{prvTrue}{$\ltrue$,}{}
$\Atom{R}{t_1, \dots, t_n}$,
$\eq[t_1][t_2]$\\
$\lnot$ & negação & $\lnot !A$ \\
$\land$ & conjunção & $(!A \land !B$) \\
$\lor$ & disjunção & $(!A \lor !B$) \\
$\lif$ & !!{conditional} & $(!A \lif !B$) \\
$\liff$ & !!{biconditional} & $(!A \liff !B)$ \\
$\lforall[][]$ & (!!{formula}) universal & $\lforall[x][!A]$ \\
$\lexists[][]$ & (!!{formula}) existencial & $\lexists[x][!A]$
\end{tabular}
\caption{Operador principal e nomes de !!{formula}s}
\ollabel{tab:main-op}
\end{table}

\end{document}
